\begin{figure*}[t]
\centering
\includegraphics[width=\linewidth]{photos/rollout.png}
\caption{\textbf{\methodname{}'s Dynamics Adaptive Behavior in} \taskone~(\textcolor{Orange}{\textit{Top}}) and \tasktwo~(\textcolor{OliveGreen}{\textit{Bottom}}) tasks. \textcolor{Orange}{\textit{Top Left:}} \methodname{} hangs the rope over the top of the container before lowering it. \textcolor{Orange}{\textit{Top Right:}} \methodname{} flips the gripper upside down before performing a direct adapter insertion. \textcolor{OliveGreen}{\textit{Bottom Left:}} \methodname{} moves the robot arm sideways and sweeps a towel horizontally from right to center to cover the bowl. \textcolor{OliveGreen}{\textit{Bottom Right}}: \methodname{} moves the arm vertically and places a hat directly on top of the bowl. This shows \methodname{}'s ability to infer different object dynamics and adapt actions accordingly to achieve the deformable object mobile manipulation tasks.}
\vspace{-15px}
\label{fig:rollout}
\end{figure*}
\section{\methodname{}: $\underlinedname$}
As shown in Fig.~\ref{fig:rda}, \methodname{} is a two-phase versatile method for learning to manipulate deformable objects in the real world by adapting to their unknown dynamics using raw visual inputs. In the first phase, \methodname{} learns a visuomotor mobile manipulation policy conditioned on two dynamics encoders using privileged ground-truth information obtained from simulation, such as ground-truth particle positions and position deltas of the deformable object  (Fig.~\ref{fig:rda}, \textit{top}). In the second phase, \methodname{} replaces the shape and dynamics encoders with two respective adaptation modules that only accept recent depth images and robot actions as input, making the entire algorithm free of privileged simulator information (Fig.~\ref{fig:rda}, \textit{middle}). At test time, \methodname{} freezes both the visuomotor policy and the two adaptation modules and deploys them directly to the real world (Fig.~\ref{fig:rda}, \textit{bottom}).

\methodname{} addresses both training phases as reinforcement learning (RL) problems which are modeled as Markov decision processes defined by the tuple $\mathcal{M} = \langle \mathcal{S}, \mathcal{O}, \mathcal{A}, \mathbb{T}, \mathcal{R}, \gamma, \rho_0, \mathcal{H}\rangle$. $\mathcal{S}$ is the state space (observable only in simulation). $\mathcal{O}$ is the space of observations. $\mathcal{A}$ is the action space. $\mathbb{T}$ is the dynamics model that governs the state transitions, $\mathbb{T}: \mathcal{S} \times \mathcal{A} \to \mathcal{S}$. $\mathcal{R}$ is a reward function. $\gamma \in [0, 1)$ is a discount factor. $\mathcal{H}$ is a finite horizon. $\rho_0$ is the distribution of initial states.
Reinforcement learning aims to train a policy, $\pi$, that takes optimal actions to maximize the expected cumulative future rewards across horizon $\mathcal{H}$.
Below, we describe each training and test phase in detail.

\textbf{Train Phase I: Learning a visuomotor policy in simulation conditioned on two ground-truth encoders}.
In the first phase, \methodname{} trains a visuomotor policy, $\pi$, that generates full-DOF robot actions to achieve the deformable object mobile manipulation task in simulation (Fig.~\ref{fig:rda}, \textit{top}). Concretely, the policy accepts a single depth image $o_t$ and a Dynamics Embedding $z_t^{d}$ as input. The Dynamics Embedding is encoded by a Dynamics Encoder $\mu_{d}$, which accepts the masses and positions of objects in simulation, and a Shape Embedding $z_t^{s}$ as input. As such, this Dynamics Encoder encodes dynamics information crucial to achieving the task. The Shape Embedding is, in turn, encoded by a Shape Encoder $\mu_{s}$, which accepts ground-truth object shape information as input, such as recent positions of the deformable object's particles and corresponding robot actions. As such, this Shape Encoder encodes information about how the shape of the object changes, which is crucial for achieving the task. We learn to encode dynamics and shape information in two separate encoders to quantify the dynamics vs. shape adaptation module's training losses we will encounter in the second phase. Because the dynamics and shape information are represented by low-dimensional privileged information from the simulator instead of high-dimensional visual observations, we train both the policy and the dynamics and shape encoders end-to-end using RL. 

Once \methodname{} has trained the visuomotor policy and the shape and dynamics encoders end-to-end, it learns two new adaptation modules --a shape adaptation module and a dynamics adaptation module--- to replace the shape and dynamics encoders, respectively, which we discuss next. 

\textbf{Train Phase II: Learning shape and dynamics adaptation modules in simulation}.
For \methodname{} to work in the real world, its policy, shape encoder, and dynamics encoder cannot accept any privileged ground-truth simulator information as input. Therefore, in this second training phase, the shape and dynamics encoders are replaced with their respective shape and dynamics adaptation modules that can be trained to infer the same dynamics embedding using only visual observations (Fig.~\ref{fig:rda}, \textit{middle}). Concretely, we first replace the Shape Encoder with a new Shape Adaptation module, $\phi_{s}$, which accepts recent depth images and robot joint angles and actions instead of ground-truth particle positions of the object as input. To train this Shape Adaptation module, we generate the ground-truth Shape Embedding, $z_t^{s}$, using the Shape Encoder and the inferred Shape Embedding, $\hat{z}_t^{s}$, from the Shape Adaptation module, and use L1-Loss to regress the inferred embedding to the ground-truth embedding. 
Similarly, we replace the Dynamics Encoder by training a new Dynamics Adaptation module, $\phi_{d}$, which accepts recent depth images and robot joint angles and actions instead of ground-truth dynamics information as input, using L1-loss between the ground-truth Dynamics Embedding $z_t^{d}$ and the inferred Dynamics Embedding $\hat{z}_t^{d}$. To prevent the Shape Adaptation module from encoding redundant dynamics information, we stop upstream gradients of training the Dynamics Adaptation module from flowing into the Shape Adaptation module. Finally, we fine-tune the visuomotor policy obtained during the first phase using the two new embeddings generated by the two new adaptation modules using RL. We stop upstream RL gradients from flowing into both adaptation modules so that the two modules separately encode the object dynamics and shape information-- even though they observe the same input-- that cannot be directly observed by the depth image $o_t$. The dynamics and shape encoders from Phase I are discarded before test time deployment, as discuss below.

\textbf{Test Time Deployment}. At test time, the robot is randomly placed in the real world with previously unseen deformable and rigid-body objects arranged randomly on a table. It needs to perform a task that requires using a deformable object with unknown dynamics. For each timestep, the robot interacts with objects in the environment using \methodname{}'s visuomotor policy, which accepts the current robot observation (i.e., a depth image and joint angles) and the Dynamics Embedding $\hat{z}_t^{d}$ as input and outputs full-DOF robot actions (Fig.~\ref{fig:rda}, \textit{bottom}). To generate the Dynamics Embedding, the robot stores a running buffer of the 10 most recent depth-image-observation-action pairs as input into both the shape and dynamics adaptation modules. To maintain short-term temporal behavioral consistency of the visuomotor policy, \methodname{} only periodically updates the shape and dynamics embeddings every 5 timesteps. This cycle repeats until the robot completes the task successfully. 

